\subsubsection{Testcase module 1: Lập kế hoạch AI}

\hypertarget{target:tc_m1}{}
\begin{landscape}
\small 
\begin{longtable}{|p{2.5cm}|p{4.0cm}|p{6.5cm}|p{4.5cm}|p{4.5cm}|}
    
    % --- PHẦN 1: TIÊU ĐỀ CHO TRANG ĐẦU TIÊN ---
\caption{Danh sách Test Case chi tiết - Module 1: Lập kế hoạch AI} \label{tab:tc_module1} \\
\hline
\textbf{ID} & \textbf{Tên Test Case} & \textbf{Các bước thực hiện} & \textbf{Dữ liệu kiểm thử} & \textbf{Kết quả mong đợi} \\
\hline
\endfirsthead

% --- PHẦN 2: TIÊU ĐỀ CHO CÁC TRANG TIẾP THEO ---
\multicolumn{5}{c}{{\bfseries \tablename\ thetable{}: Danh sách Test Case chi tiết - Module 1: Lập kế hoạch AI}} \\
\hline
\textbf{ID} & \textbf{Tên Test Case} & \textbf{Các bước thực hiện} & \textbf{Dữ liệu kiểm thử} & \textbf{Kết quả mong đợi} \\
\hline
\endhead

\hline
\endlastfoot

% ================= NỘI DUNG BẢNG =================

% --- LUỒNG CHÍNH ---
TC\_M1\_01 & Tạo lộ trình thành công (Full Option) & 
1. Đăng nhập. \newline 
2. Nhấn ``Tạo lộ trình mới''. \newline 
3. Nhập đầy đủ các trường bắt buộc và tùy chọn (Ngân sách, Sở thích). \newline 
4. Nhấn ``Lập kế hoạch''. & 
- Điểm đến: Đà Lạt \newline 
- Ngày: 01/12 - 03/12 \newline 
- Ngân sách: Trung bình \newline 
- Sở thích: Cà phê, Check-in & 
Hệ thống hiển thị bản lộ trình nháp (Timeline view) với các địa điểm phù hợp tiêu chí đã chọn. \\
\hline

TC\_M1\_08 & Đề xuất dịch vụ thành công & 
1. Mở chi tiết lộ trình. \newline 
2. Chọn địa điểm ``Địa đạo Củ Chi''. \newline 
3. Nhấn ``Gợi ý dịch vụ/Tour''. & 
- Địa điểm: Củ Chi (Tags: Di tích, Xa trung tâm) & 
Hệ thống hiển thị danh sách Tour trọn gói hoặc vé tham quan tương ứng từ đối tác. \\
\hline

TC\_M1\_13 & Lưu lộ trình MỚI thành công & 
1. Nhấn ``Lưu lộ trình'' (Lộ trình chưa từng lưu). \newline 
2. Nhập tên chuyến đi hợp lệ. \newline 
3. Nhấn Xác nhận. & 
- Tên: ``Hè Nha Trang 2025'' & 
Hệ thống thông báo ``Đã lưu thành công'' và dữ liệu được ghi vào CSDL. \\
\hline

TC\_M1\_14 & Lưu đè lộ trình CŨ thành công & 
1. Mở lộ trình đã lưu. \newline 
2. Chỉnh sửa nội dung. \newline 
3. Nhấn ``Lưu lộ trình''. & 
N/A & 
Hệ thống lưu trực tiếp các thay đổi vào CSDL, không hiển thị popup nhập tên. \\
\hline

TC\_M1\_16 & Xóa lộ trình thành công & 
1. Tại danh sách lộ trình, chọn 1 dòng. \newline 
2. Nhấn icon Thùng rác. \newline 
3. Popup hiện ra, nhấn ``Đồng ý''. & 
N/A & 
Lộ trình biến mất khỏi danh sách hiển thị và bị xóa khỏi CSDL. Thông báo thành công. \\
\hline

% --- LUỒNG THAY THẾ ---
TC\_M1\_02 & Tạo lộ trình giá trị mặc định & 
1. Nhấn ``Tạo lộ trình mới''. \newline 
2. Chỉ nhập Điểm đến, Thời gian. \newline 
3. Bỏ qua Ngân sách, Sở thích. \newline 
4. Nhấn ``Lập kế hoạch''. & 
- Điểm đến: TP.HCM \newline 
- Ngày: 2 ngày & 
Hệ thống tự động áp dụng giá trị mặc định (Ngân sách TB, Sở thích phổ thông) và tạo lộ trình thành công. \\
\hline

TC\_M1\_03 & Tạo lộ trình bằng Nhân bản (Clone) & 
1. Vào ``Lịch sử chuyến đi''. \newline 
2. Chọn lộ trình cũ, nhấn ``Tạo chuyến đi từ lộ trình này''. \newline 
3. Nhập Ngày khởi hành mới. \newline 
4. Nhấn ``Xác nhận''. & 
- Ngày cũ: 01/01/2023 \newline 
- Ngày mới: 10/12/2025 & 
Hệ thống sao chép danh sách địa điểm. Ngày giờ các hoạt động được tính toán lại theo ngày khởi hành mới. \\
\hline

TC\_M1\_11 & Tối ưu hóa có sử dụng tính năng ``Khóa'' & 
1. Chọn ``Tối ưu hóa bằng AI''. \newline 
2. Chọn ``Khóa'' (Pin) địa điểm B. \newline 
3. Chọn tùy chọn ``Tối ưu đường đi''. \newline 
4. Nhấn Áp dụng. & 
- Địa điểm B: Quán ăn ngon (đã Pin) & 
Hệ thống trả về lộ trình mới: Địa điểm B vẫn giữ nguyên vị trí/khung giờ, các địa điểm khác được sắp xếp lại. \\
\hline

TC\_M1\_12 & Thay thế địa điểm (Smart Swap) & 
1. Chọn địa điểm C (không thích). \newline 
2. Nhấn ``Gợi ý thay thế''. \newline 
3. Chọn địa điểm D từ danh sách gợi ý. & 
- Địa điểm C: Bảo tàng \newline 
- Địa điểm D: Triển lãm tranh & 
Địa điểm D thay thế vị trí của C trong lộ trình, thời gian và chi phí toàn chuyến được cập nhật lại. \\
\hline

TC\_M1\_17 & Hủy thao tác xóa lộ trình & 
1. Tại danh sách lộ trình, chọn 1 dòng. \newline 
2. Nhấn icon Thùng rác. \newline 
3. Popup hiện ra, nhấn ``Hủy''. & 
N/A & 
Hộp thoại đóng lại, lộ trình vẫn còn nguyên trong danh sách. \\
\hline

TC\_M1\_06 & AI Check Thời tiết (Tích hợp API) & 
1. Tạo lộ trình mới. \newline 
2. Hệ thống tự động gọi API thời tiết cho các ngày đã chọn. & 
- Mock API Response: \{ ``forecast'': ``Heavy Rain'' \} & 
Hiển thị cảnh báo (Alert) trên timeline: ``Dự báo có mưa lớn, hãy chuẩn bị phương án trong nhà.'' \\
\hline

% --- TEST CASE BỔ SUNG ---
TC\_M1\_18 & Đổi tên lộ trình (Rename) & 
1. Mở chi tiết lộ trình. \newline 
2. Nhấn vào tiêu đề lộ trình. \newline 
3. Nhập tên mới và nhấn Enter (hoặc click ra ngoài). & 
- Tên cũ: ``Hè Nha Trang'' \newline
- Tên mới: ``Summer Trip 2025'' & 
Tên lộ trình được cập nhật thành công trên giao diện và CSDL mà không cần tải lại trang. \\
\hline

TC\_M1\_19 & Xuất lộ trình ra file (Export) & 
1. Mở chi tiết lộ trình. \newline 
2. Nhấn nút ``Chia sẻ/Xuất''. \newline 
3. Chọn định dạng PDF. & 
- Định dạng: PDF & 
Hệ thống tải xuống file PDF chứa thông tin chi tiết lịch trình, bản đồ và danh sách chi phí dự kiến. \\
\hline

% --- LUỒNG NGOẠI LỆ ---
TC\_M1\_04 & Kiểm tra lỗi logic thời gian & 
1. Tại form tạo lộ trình. \newline 
2. Nhập Ngày bắt đầu $>$ Ngày kết thúc. \newline 
3. Nhấn ``Lập kế hoạch''. & 
- Start: 05/12/2025 \newline 
- End: 01/12/2025 & 
Hệ thống hiển thị thông báo lỗi đỏ ngay tại trường nhập liệu: ``Ngày kết thúc không thể trước ngày bắt đầu''. \\
\hline

TC\_M1\_20 & Kiểm tra dữ liệu bắt buộc (Trống) & 
1. Nhấn ``Tạo lộ trình mới''. \newline 
2. Bỏ trống trường ``Điểm đến''. \newline 
3. Nhấn ``Lập kế hoạch''. & 
- Điểm đến: [Trống] \newline
- Ngày: Hợp lệ & 
Nút ``Lập kế hoạch'' bị vô hiệu hóa (disabled) hoặc hiển thị thông báo lỗi ``Vui lòng nhập điểm đến''. \\
\hline

TC\_M1\_05 & AI không tìm thấy lộ trình & 
1. Nhập tiêu chí mâu thuẫn (Ngân sách thấp cho nơi sang trọng). \newline 
2. Nhấn ``Lập kế hoạch''. & 
- Ngân sách: 50k VND \newline 
- Điểm đến: Resort 5 sao & 
Hệ thống thông báo: ``Không thể tạo lộ trình phù hợp. Vui lòng nới lỏng yêu cầu về ngân sách hoặc địa điểm''. \\
\hline

TC\_M1\_07 & API Thời tiết lỗi/Mất mạng & 
1. Tạo lộ trình mới. \newline 
2. Giả lập mất kết nối Internet hoặc API thời tiết chết. & 
- Kết nối: Timeout & 
Lộ trình vẫn được tạo thành công nhưng không có thẻ cảnh báo thời tiết. Hệ thống không được crash. \\
\hline

TC\_M1\_09 & Đề xuất dịch vụ (Không dữ liệu) & 
1. Chọn một địa điểm chưa có đối tác liên kết. \newline 
2. Nhấn ``Gợi ý dịch vụ''. & 
- Địa điểm: Công viên ghế đá tự phát & 
Hệ thống thông báo (Toast/Popup): ``Hiện chưa có dịch vụ liên kết cho địa điểm này''. \\
\hline

TC\_M1\_10 & Cảnh báo xung đột giờ mở cửa & 
1. Kéo địa điểm A sang khung giờ 18:00. \newline 
2. Địa điểm A đóng cửa lúc 17:00. & 
- Địa điểm A: Bảo tàng (Mở 8:00 - 17:00) & 
Hệ thống hiển thị cờ cảnh báo đỏ (Warning flag): ``Địa điểm này đóng cửa lúc 17:00'' nhưng vẫn cho phép lưu nếu người dùng muốn. \\
\hline

TC\_M1\_15 & Lỗi trùng tên khi lưu mới & 
1. Nhấn ``Lưu lộ trình''. \newline 
2. Nhập tên đã tồn tại trong danh sách của tôi. \newline 
3. Nhấn Xác nhận. & 
- Tên: ``Hè Nha Trang 2025'' (đã có) & 
Hệ thống báo lỗi: ``Tên lộ trình đã tồn tại, vui lòng chọn tên khác'' và không cho phép lưu. \\
\hline

\end{longtable}
\end{landscape}